\section{The marginal coefficients from finite insertions}\label{sec:coefficients}
Linear transport leaves the variance density unchanged.  Its first nonlinear
increment, and the first increments of marked charge moments, must therefore
be computed.  We compute their averages over scales by comparing linear
charges of logarithms with normalized charges of finite products of
microscopic insertions.  The latter charges are preserved by the exact
partition identities.  Only coefficients through degree four are used;
there is no infinite-volume random power series in this argument.

\subsection{A bounded-degree test and its profiles}
Give every microscopic site $i$ a formal variable $t_i$ and start with the
factor $1+t_i e(i)$.  Initially its logarithm is put into the principal slot
of its cell.  First work with finitely many sites and retain only the degrees
needed below.  Apply the ordinary Taylor rules at the zero background, and
then substitute
\begin{equation}\label{coef:test}
 t_i=\lambda\eta_i,\qquad
 \Prob(\eta_i=1)=\Prob(\eta_i=-1)=\tfrac12,
\end{equation}
with independent signs.  Coefficient notation $[\lambda^k]$ means the
coefficient, without a factorial.  At fixed degree and fixed scale every
coefficient has bounded range.  A large finite window may consequently tend
to the whole plane after that coefficient has been computed.  Averaging over
the signs gives homogeneous arrays in the norms of Section~\ref{sec:linear}.
For a mark at $y$ start with $r_0^d\sigma_y$ and use scaled marked transport.

Let $h_j$ be the formal background list and $u_{j,y}$ its scaled marked
list.  Define
\begin{align*}
 m_j^*&=[\lambda^2]\E h_j,&
 K_j^*&=[\lambda^2]\Cov(h_j,h_j),\\
 W_j^*&=[\lambda^4]\cum(h_j,h_j,h_j),&
 Y_j^*&=[\lambda^4]\cum(h_j,h_j,h_j,h_j).
\end{align*}
For a fixed positive integer $q$ the $q$ source indices remain separate in
$u_{j,y}^{\otimes q}$, but that tensor is one grouped random argument in the
joint cumulant
\begin{equation}\label{coef:Z}
 Z_{j,q,y}^*=[\lambda^2]\cum(u_{j,y}^{\otimes q},h_j).
\end{equation}
Put
\begin{align}
 \beta_j&=[\lambda^4]g_{r_j}\bigl(\Cov(h_j,h_j)\bigr),
       &b_j&=\beta_{j+1}-\beta_j,\label{coef:b}\\
 \alpha_{j,q}(y)&=[\lambda^2]\E
                   Q_{j,q,y}(u_{j,y}^{\otimes q}),
       &c_{j,q}(y)&=\alpha_{j+1,q}(y)-\alpha_{j,q}(y).
       \label{coef:c}
\end{align}
The definitions include any bounded initial-scale discrepancy between a
microscopic list and the averaged representatives of Section~\ref{sec:linear}.

\begin{lemma}[Uniform formal profiles]\label{coef:profiles}
For $j\ge1$, the arrays $m_j^*,W_j^*,Y_j^*,Z_{j,q,y}^*$ have bounded indexed
projective norms, uniformly in $j$ and $y$ and for every fixed finite set of
$q$'s.  Moreover $K_j^*=V_j$ and the degree-zero marked list is $w_{j,y}$.
The coefficients $b_j$ and $c_{j,q}(y)$ are uniformly bounded.
\end{lemma}

\begin{proof}
Linear insertions propagate by nested averaging.  Thus the coefficient of
$t_i$ at scale $r$ is $P_{Hr,c(i)}e(i)$ in the principal slot of its current
cell, giving $K_j^*=V_j$ exactly, and the same statement for the spin gives
$w_{j,y}$.  We explain separately why the mean profile stays bounded despite
the expanding mean direction.

For one microscopic variable, write $H_{i,k}^{(j)}$ for its $t_i^k$ coefficient
in the plane-centered background log list at scale $j$.  All its anchors
follow the same chain of parent cells.  The fixed-degree carrier estimate
of Lemma~\ref{block:carrier} bounds the number of occupied slot types
and their lengths by constants independent of $j$.  Increasing $L$ makes
these slots short for linear transport.  Constants in the present formal
calculation may depend on $L$ and $r_0$.  We claim
\begin{equation}\label{coef:single-columns}
 \|H_{i,1}^{(j)}\|\le Cr_j^{-1},\qquad
 \|H_{i,2}^{(j)}\|\le Cr_j^{-2}.
\end{equation}
The first bound is the reference energy estimate.  At a step from radius
$s$ to $Ls$, every freshly generated quadratic term costs $C\|H_{i,1}\|^2$,
hence at most $C_L(Ls)^{-2}$.  To bound the transported old quadratic terms,
condition on one contour outside all earlier cuts.  The exact identity,
with saved constants restored, represents the original factor $1+t_i e(i)$.
Its quadratic coefficient is zero.  Its first scalar coefficient is zero
because $e(i)$ is plane centered.  Equating quadratic coefficients therefore
gives, modulo constants,
\begin{equation}\label{coef:single-cancel}
 P H_{i,2}=-\tfrac12 P(H_{i,1}^2).
\end{equation}
Here each $H$ denotes the sum of its list entries before this common outer
average.  The linear rule puts these old short entries in the same parent
principal slot; cancellation in~\eqref{coef:single-cancel} thus occurs
before taking its norm.  The right side is bounded by $C s^{-2}$.
Combining old and fresh terms proves the second bound at the next scale.
At the first scale the bound is included by the allowed constant depending
on $L,r_0$.  This reasoning also covers earlier saved scalar coefficients:
they are reinstated for the exact identity and disappear upon plane
centering.  Their footprints obey the same carrier bound.

Only squares of individual signs contribute to $[\lambda^2]\E h_j$.
There are at most $C r_j^2$ microscopic sites which can contribute to a fixed
output anchor.  Summing the second estimate in~\eqref{coef:single-columns}
proves the bound on $m_j^*$.  The finitely many fixed-degree slot lengths
permit the same conclusion with the indexed exponential weights.

For completeness, the remaining profile recurrences follow by the
moment--cumulant partition formula.  The lowest nonzero fourth-cumulant
coefficient has four linear background letters, so $Y_j^*$ transports
homogeneously by $T_j^{(4)}$.  At degree four the third cumulant has, besides
its homogeneous term, only a quadratic output letter and two linear output
letters.  Its input partitions use either $Y_j^*$ or two copies of $V_j$.
These terms are bounded by the multilinear estimates of
Proposition~\ref{norm:taylor}.  The strict contractions of
Proposition~\ref{lin:background} therefore bound $Y_j^*$ and $W_j^*$.
In~\eqref{coef:Z}, the constant marked tensor is deterministic.  Its first
nonzero covariance with a background uses one linear marked correction and
one linear background letter.  The recurrence is the contracting connected
marked/background transport, with bounded forcing from the constant marked
tensor and $V_j$; hence $Z_{j,q,y}^*$ is bounded by
Proposition~\ref{lin:marked}.  The degree-two coefficient of
$\cum(u_{j,y}^{\otimes q},h_j,h_j)$ is zero: at that degree the marked
argument is deterministic, and a joint cumulant with a deterministic
argument vanishes.

Finally, in the degree-four covariance increment the linearly transported
covariance has identical $g$ coordinate.  The remaining partitions use
$W_j^*,Y_j^*,V_j\otimes V_j$, or $m_j^*\otimes V_j$, all bounded.  In the
degree-two marked-charge increment the transported marked tensor likewise
has unchanged charge; the remaining terms use $m_j^*,V_j$ and
$Z_{j,q,y}^*$, while the connected two-background profile just discussed
vanishes.  This proves the asserted bounds on the increments.
\end{proof}

\subsection{Charges invariant under a finite sequence of cuts}
For finitely many microscopic insertions let $H_*^{(j)}$ be the sum of the
plane-centered log-list entries produced through scale $j$.  At any fixed
coefficient order this is a finite sum with its full carrier footprints.
Let $f_i^{(j)}$ be the unscaled coefficient of one square-free spin source
in log-list slot $i$, and set $f^{(j)}=\sum_i f_i^{(j)}$.  The scaled
source entries are $u_{j,i}=r_j^d f_i^{(j)}$.  In particular,
\[
 Q_{j,q,y}(u_j^{\otimes q})=\ell_s(f^{(j)})^q.
\]
This identity connects the scalar charge calculation below to the indexed
tensor coordinate in~\eqref{coef:c}.  Define formal normalized charges by
\begin{equation}\label{coef:invariants}
 \mathcal E_j=\frac{\ell(\exp H_*^{(j)})}
                    {\langle\exp H_*^{(j)}\rangle_0},\qquad
 \mathcal S_j=\frac{\ell_s(f^{(j)}\exp H_*^{(j)})}
                    {\langle\exp H_*^{(j)}\rangle_0}.
\end{equation}
Both denominators have constant term one, so their formal reciprocals are
well defined.  Saved background constants cancel between numerator and
denominator.  A saved odd scalar is zero by spin-flip symmetry.

\begin{lemma}[Finite-insertion invariance]\label{coef:invariance}
Each coefficient of $\mathcal E_j$ and $\mathcal S_j$ in the formal variables
is independent of $j$.  It equals the corresponding normalized charge of
the original microscopic product and, for $\mathcal S_j$, its spin insertion.
\end{lemma}

\begin{proof}
Fix the coefficient and place an outer contour beyond every cut and carrier
which occurs in it.  For each fixed spin assignment on that contour, apply
Proposition~\ref{block:exact} to the system inside it.  The termwise
conditional-expectation proof of that proposition leaves these contour spins
fixed: each inner cut is an application of the tower property within this
conditional reference measure.  Thus the conditional partition expression,
as a function of the outer spin assignment, agrees before and after every
step.  Its source derivative gives the analogous identity with the single
odd insertion.  Let $C_j$ be the saved scalar formal series, so that the restored
weight is $e^{C_j}\exp H_*^{(j)}$.  Taking the reference plane average
proves that $e^{C_j}\langle\exp H_*^{(j)}\rangle_0$ is independent
of $j$.  Taking the disk-response limit defining $\ell$ or $\ell_s$
proves the corresponding invariance of each charge numerator multiplied
by $e^{C_j}$.  The saved scalar with one odd source is zero, so the marked
identity has this same scalar factor.  Every operation involves a finite
coefficient and finitely many cuts; no interchange with a convergent
infinite formal expansion is involved.  Dividing cancels $e^{C_j}$ and
proves the claim.
\end{proof}

We next quantify the difference between these invariant charges and the
linear charges $\ell(H_*^{(j)})$ and $\ell_s(f^{(j)})$ which occur in the
Taylor coordinates.  In the following bounds set $r=r_j$ and
$h=1+|i-k|$ for two distinct sites.  For $h\le Cr$ we have
\begin{equation}\label{coef:mixed-column}
 \|[t_it_k]H_*^{(j)}\|\le \frac{C}{hr}.
\end{equation}
The norm here is the sum of the finitely many occupied slot norms, with
fixed-degree weights; it also bounds their aggregate sup norm.  A mixed
coefficient cannot be generated before the two microscopic descendants have
contacting carriers.  This requires a generation radius $s\ge c_Lh$.
At radius $s$ a new quadratic coefficient costs $C_Ls^{-2}$.
Subsequent linear transport along ancestor cells is nested conditional
averaging and has centered-even bound $C_Ls/r$.  Anchor splitting adds only
a fixed-degree number of terms, already included in the constant.  Thus
\[
 \sum_{\substack{s=r_b\le r\\s\ge c_Lh}}
       C_Ls^{-2}\frac{s}{r}
 \le \frac{C_L}{r}\sum_{b:r_b\ge c_Lh}r_b^{-1}
 \le \frac{C_L}{hr}.
\]
This also covers $h$ below the starting radius by changing the constant.
The last generation has no further averaging and obeys the same bound.

For one energy site and one spin mark, put $h=1+|i-y|$.  The corresponding
bounds are
\begin{equation}\label{coef:marked-columns}
 \|[1]f^{(j)}\|\le Cr^{-d},\qquad
 \|[t_i]f^{(j)}\|\le Cr^{-d}h^{-1}.
\end{equation}
Indeed a mixed spin/energy generation at radius $s\ge c_Lh$ costs
$C_Ls^{-1-d}$ and subsequent odd averaging costs $C_L(s/r)^d$.
The geometric sum is $C_Lr^{-d}\sum_{b:r_b\ge c_Lh}r_b^{-1}$, as claimed.
All supports in these comparisons fit in a square of radius $C_Lr$.
Charge extraction therefore costs at most $C_Lr$ for even functions and
$C_Lr^d$ for odd functions.

\subsection{The covariance increment}
\begin{proposition}[Averages of the marginal coefficients]\label{coef:averages}
The coefficients in~\eqref{coef:b}--\eqref{coef:c} satisfy
\begin{align}
 \sum_{j<J}b_j&=-4\pi a^2\log r_J+o(\log r_J),\label{coef:bsum}\\
 \sum_{j<J}c_{j,q}(y)&=\binom q2\frac{\pi a^2}{2}\log r_J
                         +o(\log r_J),\label{coef:csum}
\end{align}
where the second error is uniform in $y$ for each fixed $q$.  Equivalently,
their Ces\`aro means are $-4\pi a^2\log L$ and
$\binom q2(\pi a^2/2)\log L$.
\end{proposition}

\begin{proof}
We first prove~\eqref{coef:bsum}.  The coordinate $\beta_J$ is the
coefficient of degree four in the variance per unit microscopic area of
$\ell(H_*^{(J)})$.  To see this directly, restrict to a large union of cells,
sum its linear charges, divide its variance by the number of microscopic
sites, and let the window grow.  At fixed $J$ the coefficient arrays have
bounded range, so the boundary contribution vanishes.  Formula~\eqref{lin:g}
is precisely the resulting anchored covariance sum.  At degree four,
averaging the independent signs leaves only one-site terms and unordered
two-site terms.  Terms supported on disjoint noninteracting insertions
cancel in the variance subtraction.  The surviving pair separations are
at most $C_Lr_J$, by the fixed-degree carrier bound.

Fix two sites $i\ne k$ and abbreviate
\[
 C_{ik}=\langle e(i)e(k)\rangle_0,\qquad
 A_{ik}=\ell(e(i)e(k)).
\]
Their invariant energy charge is exactly
\begin{equation}\label{coef:pair-rational}
 \mathcal E=
 \frac{t_i+t_k+A_{ik}t_it_k}{1+C_{ik}t_it_k}.
\end{equation}
The numerator uses $\ell(e(i))=\ell(e(k))=1$ and $\ell(1)=0$;
the denominator uses the centering of each energy.  Through degree three,
\begin{equation}\label{coef:pair-expansion}
 \mathcal E=t_i+t_k+A_{ik}t_it_k
              -C_{ik}(t_i^2t_k+t_it_k^2)+O(t^4).
\end{equation}
Here $O(t^4)$ denotes only discarded formal monomials of total degree at
least four.  Consequently the degree-four variance contribution after
$t_i=\lambda\eta_i$, $t_k=\lambda\eta_k$ is
\begin{equation}\label{coef:pair-value}
 A_{ik}^2-4C_{ik}.
\end{equation}
The square of the quadratic mixed term supplies $A_{ik}^2$; its mean is
zero.  The two linear/cubic crosses each supply $-2C_{ik}$.
Pure quadratic terms, if present in a different log representation, are
deterministic after the sign substitution and disappear in the variance.

We justify replacing the linear charge by this invariant charge with a
bounded total error.  Write $A_i=[t_i]H_*$, $B_i=[t_i^2]H_*$ and
$D_{ik}=[t_it_k]H_*$ at scale $r=r_J$.  The bounds above give
\[
 \|A_i\|\le C/r,\quad \|B_i\|\le C/r^2,\quad
 \|D_{ik}\|\le C/(hr).
\]
All columns are plane centered.  In particular the denominator in
\eqref{coef:invariants} has zero first-degree term.  At mixed degree two,
$\exp H_*-H_*$ contributes $A_iA_k$, whose energy charge is $O(r^{-1})$.
Thus the linear charge has first-degree coefficients exactly one and
mixed second-degree coefficient $A_{ik}+O(r^{-1})$.
At $t_i^2t_k$, the numerator corrections beyond the linear log charge
are the charges of
\[
 A_iD_{ik},\qquad B_iA_k,\qquad \tfrac12 A_i^2A_k.
\]
The denominator correction is its quadratic coefficient multiplied by a
first-degree numerator.  That coefficient is bounded by
$C(r^{-2}+(hr)^{-1})$ (in fact the plane-centering cancellation may improve
some of its terms).  Multiplying the displayed sup-norm products by the
charge cost $Cr$ gives the same bound.  Therefore the linear and invariant
coefficients at $t_i^2t_k$ and $t_it_k^2$ differ by
\begin{equation}\label{coef:cubic-error}
 O\bigl(r^{-2}+(hr)^{-1}\bigr).
\end{equation}
Since $|A_{ik}|\le C/h$ by Lemma~\ref{ref:ope}, squaring the
second-degree coefficient adds at most
$C((hr)^{-1}+r^{-2})$.  The first-degree/cubic crosses obey the same bound.
This proves that the pair contribution differs from~\eqref{coef:pair-value}
by at most $C((hr)^{-1}+r^{-2})$.

For a single site the invariant charge is exactly $t_i$.  The pure cubic
charge discrepancy consists of the products $A_iB_i$, $A_i^3/6$, and the
quadratic denominator term times the linear numerator, and is $O(r^{-2})$.
Hence its degree-four variance contribution is $O(r^{-2})$ per site.
No bound on the third log column itself is being assumed: its linear charge
is determined by the invariant identity and these lower-column corrections.

The pair errors have a bounded sum per unit area, since in two dimensions
\begin{equation}\label{coef:error-sum}
 \sum_{1\le |z|\le C_Lr}
        \left(\frac1{r(1+|z|)}+\frac1{r^2}\right)\le C_L.
\end{equation}
Indeed the shell $n\le |z|<n+1$ has at most $C(n+1)$ lattice points; the
first summand contributes $C/r$ per shell and the second has at most
$C r^2$ terms.  The one-site contribution is bounded as well.
One may use the disk $|z|\le C_Lr$ for the pair sum: any noninteracting
pair in this range has zero linear variance contribution and satisfies
the same comparison, and beyond this range the linear coefficient is zero.
Thus
\begin{equation}\label{coef:beta-lattice}
 \beta_J=\frac12\sum_{0<|z|\le C_Lr_J}
       \left\{\ell(e(z)e(0))^2-4\langle e(z)e(0)\rangle_0\right\}+O(1).
\end{equation}
The half accounts for unordered pairs.  If the fixed-scale list is only
grid-translation invariant, the intermediate expression averages its
linear pair contribution over the first site's residue class in a grid
cell; the invariant expression is fully translation invariant and the
uniform error bound gives exactly~\eqref{coef:beta-lattice}.

Lemma~\ref{ref:ope} gives, uniformly in direction,
$\langle e(z)e(0)\rangle_0=a^2|z|^{-2}+o(|z|^{-2})$ and
$\ell(e(z)e(0))^2=o(|z|^{-2})$.  The lattice sum satisfies
\[
 \sum_{0<|z|\le R}|z|^{-2}=2\pi\log R+O(1).
\]
For example compare each unit lattice square away from the origin with
its radial integral; the gradient error is $O(|z|^{-3})$, which is summable
outside any fixed disk.  For an $o(|z|^{-2})$ remainder, split at a fixed
large radius, make its coefficient arbitrarily small beyond that radius,
and use the same logarithmic bound.  Its sum is consequently $o(\log R)$.
Applying these facts in~\eqref{coef:beta-lattice} yields
$\beta_J=-4\pi a^2\log r_J+o(\log r_J)$.
Telescoping~\eqref{coef:b} subtracts only the bounded initial coefficient
$\beta_0$, proving~\eqref{coef:bsum}.

\subsection*{The marked increment}
For one energy site $i$ and mark $y$, the invariant normalized spin charge
is exactly
\begin{equation}\label{coef:spin-invariant}
 \mathcal S=1+t_i S_i(y),\qquad
 S_i(y)=\ell_s(\sigma_y e(i)).
\end{equation}
The denominator is one because the energy is centered.  Write
$F_0=[1]f$, $F_1=[t_i]f$, and retain the background columns $A_i,B_i$
above.  At first degree the difference between
$\ell_s(f)$ and $\mathcal S$ is $\ell_s(F_0A_i)$, which is $O(r^{-1})$
by~\eqref{coef:marked-columns}.  At second degree the numerator corrections
are the charges of
\[
 F_1A_i,\qquad F_0B_i,\qquad \tfrac12F_0A_i^2.
\]
Their sup norms, multiplied by the odd charge cost $Cr^d$, are bounded by
$C((hr)^{-1}+r^{-2})$.  The denominator correction has the same bound.
It follows that the linear marked charge is
\[
 1+t_i\{S_i(y)+O(r^{-1})\}
       +t_i^2 O\bigl((hr)^{-1}+r^{-2}\bigr)+O(t_i^3).
\]
Its $q$th power has degree-two coefficient
\begin{equation}\label{coef:spin-pair-error}
 \binom q2 S_i(y)^2+O_q\bigl((hr)^{-1}+r^{-2}\bigr),
\end{equation}
since $|S_i(y)|\le C/h$.  Sign averaging eliminates distinct-site quadratic
monomials, so only this one-site calculation is needed.  The marked
coefficient has range at most $C_Lr$; sum~\eqref{coef:spin-pair-error} over
those sites and use~\eqref{coef:error-sum}.  Uniformly in $y$,
\begin{equation}\label{coef:alpha-lattice}
 \alpha_{J,q}(y)=\binom q2
       \sum_{|i-y|\le C_Lr_J}\ell_s(\sigma_y e(i))^2+O_q(1).
\end{equation}
The finitely many coincident or adjacent sites contribute only a bounded
amount.  The uniform OPE asymptotic
$\ell_s(\sigma_y e(i))^2\sim a^2/(4|i-y|^2)$ and the preceding lattice
summation give~\eqref{coef:csum} after subtracting $\alpha_{0,q}(y)$.
Finally $\log r_J=J\log L+\log r_0$ gives both Ces\`aro assertions.
\end{proof}

\begin{corollary}[Bounded first-charge sums]\label{coef:bounded-first}
There is a constant independent of $J$ and $y$ such that
\[
 \left|\sum_{j<J}c_{j,1}(y)\right|\le C.
\]
\end{corollary}

\begin{proof}
For $q=1$ the invariant coefficient of $t_i^2$ in
\eqref{coef:spin-invariant} is exactly zero, rather than merely smaller
than its logarithmic scale.  Thus~\eqref{coef:alpha-lattice} gives
$|\alpha_{J,1}(y)|\le C$.  The initial coefficient is uniformly bounded,
and~\eqref{coef:c} telescopes.  This stronger conclusion will control the
first marked moment; a zero Ces\`aro mean alone would not suffice.
\end{proof}
